\scanpage{41}
% Pencil archive foliation: 40.
Nous allons expliciter plus bas la détermination
de $\sigma^0(X)$ et $\sigma^1(X)$.

Il y a lieu également de préciser
les opérations $\vee$ et $\lambda^i$. Pour la
première, notons que
\[
\unsure{\mathbb{Q}|\mathbb{Z}}(-i)[\mathcal{M}_i^+(X)]^\vee\subset \mathcal{M}_i^+(X)
\]
d'où on conclut de proche en
proche que $x\unsure{\rightsquigarrow|\mapsto} \check x\xi^i$ définit sur
$\sigma_i$ une involution. Mais en fait il
résultera de la théorie \unsure{projetée|projective} que
\uline{cette dernière est l'identité}, i.e.
\uline{l'opération} $\vee$ \uline{dans} \deleted{\unsure{\mathcal M}}
$\mathrm{M}_i(X)$ \deleted{est l'identité} \uline{n'est
autre que} $x\unsure{\rightsquigarrow|\mapsto}\xi^{-i}x$,
\deleted{\unsure{donc dans} $\mathrm{M}(X)$
$x\unsure{\rightsquigarrow|\mapsto} \check x=\xi^{-i}x$}.
Quant aux opérations $\lambda^i$, leur détermination
devrait se faire en termes de groupes finis
d'opérateurs (tels les $\unsure{\mathfrak{S}_p|G_p}$ \ldots).
